\documentclass[../ap_2017.tex]{subfiles}

\graphicspath{{./image/}}

\begin{document}

\setcounter{chapter}{4}
\chapter{ラプラス変換}
\section{}
\subsection{}
\begin{equation}\begin{aligned}[b]
    L[t^n] = \int_{0}^{\infty} t^n e^{-st} \dd{t}
    = \frac{1}{s^{n+1}}\int_{0}^{\infty}x^n e^{-x}\dd{x}
    = \frac{\Gamma(n+1)}{s^{n+1}}
    = \frac{n!}{s^{n+1}}
\end{aligned}\end{equation}
\subsection{}
\begin{equation}\begin{aligned}[b]
    L\qty[\dv{f(t)}{t}] = \int_{0}^{\infty}\dv{f(t)}{t}e^{-st}\dd{s}
    = \qty[f(t)e^{-st}]_0^\infty +s\int_{0}^{\infty}f(t)e^{-st}\dd{t}
    = sF(s)-f(0)
\end{aligned}\end{equation}

\subsection{}
\begin{equation}\begin{aligned}[b]
    L\qty[e^{at}f(t)] = \int_{0}^{\infty}f(t)e^{at}e^{-st}\dd{t}
    = \int_{0}^{\infty}f(t)e^{-(s-a)t}\dd{t}
    =F(s-a)
\end{aligned}\end{equation}

\section{}
与えられた微分方程式をラプラス変換すると
\begin{equation}\begin{aligned}[b]
    0 &= -\dv{s}\qty[s^2F(s)-sf(0)-f'(0)]+sF(s)-f(0)-3\dv{s}\qty[sF(s)-f(0)]+3F(s) \\
    &= -2sF(s)-s^2\dv{F(s)}{s}+f(0)+sF(s)-f(0)-3F(s)-3s\dv{F(s)}{s}+3F(s)\\
    &= -s\qty((s+3)\dv{F(s)}{s}+F(s))\\
    F(s) &= \frac{C}{s+3}
\end{aligned}\end{equation}
となる。最後の\(C\)はある定数。
これをラプラス逆変換して
\begin{equation}\begin{aligned}[b]
    f(t) = \frac{1}{2\pi i}\lim_{\varepsilon\to +0}\int_{\varepsilon-i\infty}^{\varepsilon+i\infty}\frac{Ce^{st}}{s+3}\dd{s}
    = Ce^{-3t}
\end{aligned}\end{equation}
初期条件を考えると\(C=1\)とわかるので、
\begin{equation}\begin{aligned}[b]
    f(t) = e^{-3t}
\end{aligned}\end{equation}

\section{}
\subsection{}
与えられた連立微分方程式をラプラス変換して
\begin{equation}\begin{aligned}[b]
    \begin{cases}
        sX(s)-x(0) &= -X(s)\\
        sY(s)-y(0) &= X(s) - 2Y(s)
    \end{cases}
    \quad \rightarrow\quad
    \begin{cases}
        X(s) &= \frac{a}{s+1}\\
        Y(s) &= \frac{a}{s+1}+\frac{b-a}{s+2}
    \end{cases}
    \quad \rightarrow\quad
    \begin{cases}
        x(t) &= ae^{-t}\\
        y(t) &= ae^{-t}+(b-a)e^{-2t}
    \end{cases}
\end{aligned}\end{equation}

\subsection{}
\begin{equation}\begin{aligned}[b]
    e^{-t} = \frac{x}{a}
\end{aligned}\end{equation}
を使うと
\begin{equation*}\begin{aligned}[b]
    y &= x + \frac{b-a}{a^2}x^2
\end{aligned}\end{equation*}
の放物線となる。

\subsection{}
こんな感じ
(\url{https://www.desmos.com/calculator/gaz6zahb2b})


\end{document}
